% =====================================================================
%  CARNET D'ENTRAINEMENT — FICHE 1 : Les chiffres significatifs
%  THÈME 1 — Compter les chiffres significatifs (le cas des zéros)
%  Fiche TUTO (à lire avant les exercices)
% =====================================================================
\documentclass[11pt,a4paper]{article}
\usepackage[utf8]{inputenc}
\usepackage[T1]{fontenc}
\usepackage{lmodern}
\usepackage[french]{babel}
\usepackage{amsmath,amssymb,mathtools}
\usepackage{icomma}
\usepackage{siunitx}
\sisetup{output-decimal-marker={,},per-mode=symbol}
\usepackage[version=4]{mhchem}
\usepackage{xcolor}
\usepackage{tikz}
\usepackage[most]{tcolorbox}
\usepackage{enumitem}
\usepackage{booktabs}
\usepackage{array}
\usepackage{fancyhdr}
\usepackage[a4paper,margin=2cm,headheight=26pt,headsep=8pt]{geometry}
\usetikzlibrary{arrows.meta,calc}

\definecolor{primary}{RGB}{150,98,162}
\definecolor{primarydark}{RGB}{92,52,110}
\definecolor{sigcol}{RGB}{0,90,200}
\definecolor{nosigcol}{RGB}{210,30,40}
\definecolor{reglecol}{RGB}{150,98,162}
\definecolor{methcol}{RGB}{90,90,100}
\definecolor{excol}{RGB}{0,135,90}
\definecolor{attncol}{RGB}{200,40,55}

\newcommand{\cs}{chiffre significatif}
\newcommand{\CS}{chiffres significatifs}

\pagestyle{fancy}\fancyhf{}
\fancyhead[L]{\small\textcolor{primarydark}{\textbf{Fiche 1} — Chiffres significatifs}}
\fancyhead[R]{\small\textcolor{primarydark}{\textsc{Tuto · Thème 1}}}
\fancyfoot[C]{\small\thepage}
\renewcommand{\headrulewidth}{0.8pt}
\renewcommand{\headrule}{\hbox to\headwidth{\color{primary}\leaders\hrule height \headrulewidth\hfill}}

\tcbset{boxbase/.style={enhanced,breakable,boxrule=0.9pt,arc=2.5pt,
  left=3mm,right=3mm,top=1.6mm,bottom=1.6mm,fonttitle=\bfseries,coltitle=white}}
\newtcolorbox{regle}[1][]{boxbase,colback=reglecol!7,colframe=reglecol,
  title={\textsc{Règle}\ifx\relax#1\relax\relax\else\ ---\ \textmd{\itshape #1}\fi}}
\newtcolorbox{methode}[1][]{boxbase,colback=methcol!7,colframe=methcol,sharp corners,
  title={\textsc{Méthode}\ifx\relax#1\relax\relax\else\ : \textmd{\itshape #1}\fi}}
\newtcolorbox{exemple}[1][]{boxbase,colback=excol!6,colframe=excol,
  title={\textsc{Exemple}\ifx\relax#1\relax\relax\else\ : \textmd{\itshape #1}\fi}}
\newtcolorbox{attention}[1][]{boxbase,colback=attncol!6,colframe=attncol,
  title={\textsc{Attention}\ifx\relax#1\relax\relax\else\ : \textmd{\itshape #1}\fi}}

\setlength{\parindent}{0pt}
\setlength{\parskip}{2pt}

\begin{document}

\begin{center}
{\LARGE\bfseries\color{primarydark} Thème 1 — Compter les chiffres significatifs}\\[1pt]
{\large\color{primary} Le cas délicat des zéros}
\end{center}

\vspace{1mm}

Un \textbf{\cs} (CS) est un chiffre qui \emph{porte une information sur la mesure}. Le
nombre de CS est la « carte d'identité » de la précision d'un nombre mesuré : plus il y en a,
plus la mesure est précise. Ce thème entraîne \emph{une seule} mécanique : \textbf{savoir compter}
les CS d'un nombre, en maîtrisant le rôle ambigu du chiffre~$0$.

\begin{regle}[le principe de base]
On lit le nombre \textbf{de gauche à droite}. Le comptage des CS \textbf{commence au premier chiffre
non nul}. \textbf{Tout chiffre non nul} ($1$ à $9$) est \emph{toujours} significatif. Toute la
difficulté se résume donc au traitement des zéros.
\end{regle}

\begin{regle}[les trois lois des zéros]
\begin{itemize}[leftmargin=1.5em,topsep=1pt,itemsep=1pt]
  \item \textbf{Zéros de tête} (à gauche du 1\textsuperscript{er} chiffre non nul) : \textcolor{nosigcol}{\textbf{jamais}}
        significatifs. Ils ne font que \emph{caler la virgule} (ordre de grandeur). Ex. : dans
        $\num{0,0072}$, les zéros $0{,}00$ ne comptent pas $\rightarrow$ \textbf{2}~CS.
  \item \textbf{Zéros captifs} (coincés \emph{entre} deux chiffres non nuls) : \textcolor{sigcol}{\textbf{toujours}}
        significatifs. Ex. : $\num{2005}$ $\rightarrow$ \textbf{4}~CS.
  \item \textbf{Zéros de queue \emph{après} la virgule} : \textcolor{sigcol}{\textbf{toujours}}
        significatifs (on ne les écrirait pas s'ils n'étaient pas mesurés). Ex. : $\num{15,200}$
        $\rightarrow$ \textbf{5}~CS ; $\num{60,0}$ $\rightarrow$ \textbf{3}~CS.
\end{itemize}
\end{regle}

\begin{exemple}[le nombre qui contient tous les cas : $m=\num{0,0035070}$]
\begin{center}
\begin{tikzpicture}[font=\small]
  \node[font=\Large] at (0,0)
     {$0,\textcolor{nosigcol}{00}\,\textbf{3}\,\textbf{5}\,\textcolor{sigcol}{\textbf 0}\,\textbf{7}\,\textcolor{sigcol}{\textbf 0}$};
  \draw[nosigcol,thick,->,>=Stealth] (-1.05,0.7) -- (-1.05,0.22);
  \node[nosigcol,above,align=center,font=\scriptsize] at (-1.05,0.7){zéros de tête\\ \emph{non} sig.};
  \draw[sigcol,thick,->,>=Stealth] (0.30,-0.72) -- (0.30,-0.22);
  \draw[sigcol,thick,->,>=Stealth] (1.18,-0.72) -- (1.18,-0.22);
  \node[sigcol,below,align=center,font=\scriptsize] at (0.74,-0.78){zéros captif et de queue\\ \emph{sont} sig.};
\end{tikzpicture}
\end{center}
Bilan : $3,\,5,\,0,\,7,\,0$ sont significatifs $\Rightarrow \boxed{\mathbf{5}}$~CS.
\end{exemple}

\begin{attention}[deux pièges à ne jamais oublier]
\textbf{1. CS $\neq$ décimales.} On ne compte pas les chiffres après la virgule, mais les chiffres
significatifs. $\num{0,007}$ a \emph{trois} décimales mais \textbf{un seul} CS ; $\num{42,5}$ a
\emph{une} décimale mais \textbf{trois} CS.

\smallskip
\textbf{2. Les zéros de queue d'un \emph{entier} sont ambigus.} Écrit tel quel, $\num{100}$ peut avoir
$1$, $2$ ou $3$~CS : on ne sait pas si les zéros sont mesurés ou de simple cadrage. Seule la
\textbf{notation scientifique} tranche : $\num{1e2}$ (1~CS), $\num{1,0e2}$ (2~CS), $\num{1,00e2}$ (3~CS).
\end{attention}

\begin{methode}[compter les CS en 3 réflexes]
\begin{enumerate}[leftmargin=1.6em,topsep=1pt,itemsep=1pt]
  \item Je repère le \textbf{premier chiffre non nul} : le comptage démarre ici.
  \item Je compte \textbf{tout} jusqu'au dernier chiffre écrit, en gardant les zéros captifs et les
        zéros de queue \emph{s'il y a une virgule}.
  \item Cas d'un entier à zéros de queue ($\num{2500}$) : je signale l'\textbf{ambiguïté} et je propose
        la notation scientifique adaptée.
\end{enumerate}
\end{methode}

\begin{exemple}[comptage express]
\small
\begin{center}\renewcommand{\arraystretch}{1.15}
\begin{tabular}{c c >{\raggedright\arraybackslash}m{7.2cm}}
\toprule
\textbf{Nombre} & \textbf{CS} & \textbf{Pourquoi}\\
\midrule
$\num{42,5}$      & $3$ & trois chiffres non nuls\\
$\num{0,0100}$    & $3$ & zéros de tête non sig. ; deux zéros de queue sig.\\
$\num{2005}$      & $4$ & deux zéros \emph{captifs} sig.\\
$\num{4,00e3}$    & $3$ & notation scientifique : tout ce qui est écrit est sig.\\
$\num{2030}$      & $3$ (ambigu) & zéro captif sig. ; zéro de queue d'entier ambigu\\
\bottomrule
\end{tabular}
\end{center}
\end{exemple}

\end{document}
